\pdfbookmark[2]{On X.}{nexus-remarks-x}
\begin{center}\large On X.\end{center}

Let $p$ be a positive number divisible by no square, of the form $6n\pm1$, and put
\[
S_r=\Sigma\left(\frac{s_r}p\right)
\]
where $s_r$ must run through all numbers relatively prime to $p$ that lie between $(r-1)\frac p{12}$ and $r\frac p{12}$. Denote the numbers of properly primitive forms that are not properly equivalent, for the determinants $-p$ and $-3p$, by $C_1$, $C_3$. One readily finds (compare \textsc{Dirichlet}, Recherches etc., § 11, or the note on VI and VII)
\[
C_1=2(S_1+S_2+S_3),\quad C_3=2(S_1+S_2-S_5-S_6)
\]
or
\[
C_1=S_1+S_2+S_3+S_4+S_5+S_6,\quad C_3=2(S_2+S_3+S_4+S_5)
\]
according as $p\equiv1$ or $\equiv3\pmod4$. Taking into account also that
\begin{center}\small
\begin{tabular}{@{}ccccccl@{}}
the&numbers&$s_1$ and $s_2$&coincide with&the&numbers&$2s_1$ and $p-2s_6$\\
„&„&$s_3$ and $s_4$&„&„&„&$2s_2$ and $p-2s_5$\\
„&„&$s_5$ and $s_6$&„&„&„&$2s_3$ and $p-2s_4$
\end{tabular}
\end{center}
and likewise
\begin{center}\small
\begin{tabular}{@{}ccccccl@{}}
the&numbers&$s_1,s_2,s_3$&coincide with&the&numbers&$3s_1,\ 3s_5-p,\ p-3s_4$\\
„&„&$s_4,s_5,s_6$&„&„&„&$3s_2,\ 3s_6-p,\ p-3s_3$
\end{tabular}
\end{center}
and that, in the case $p\equiv1\pmod4$, the sum $S_1+S_2+S_3+S_4+S_5+S_6=0$, one obtains four further relations, besides the two above, between the six sums $S_1,S_2\ldots S_6$, so that all of them can be determined from $C_1$ and $C_3$. In this way, when $p\equiv1\pmod4$, one obtains
\SourceLeaf{313}
\[
\begin{aligned}
S_1=S_{12}&=\tfrac14\left(1+\left(\frac3p\right)\right)C_1+\tfrac1{12}\left(1+\left(\frac2p\right)\right)C_3\\
S_2=S_{11}&=-\tfrac14\left(1+\left(\frac3p\right)\right)C_1+\tfrac1{12}\left(1+\left(\frac2p\right)\right)C_3\\
S_3=S_{10}&=\tfrac12C_1-\tfrac16\left(1+\left(\frac2p\right)\right)C_3\\
S_4=S_9&=-\tfrac12C_1-\tfrac16\left(2-\left(\frac2p\right)\right)C_3\\
S_5=S_8&=\tfrac14\left(1+\left(\frac3p\right)\right)C_1-\tfrac1{12}\left(5-\left(\frac2p\right)\right)C_3\\
S_6=S_7&=-\tfrac14\left(1+\left(\frac3p\right)\right)C_1+\tfrac1{12}\left(1+\left(\frac2p\right)\right)C_3
\end{aligned}
\]
and, when $p\equiv3\pmod4$,
\[
\begin{aligned}
S_1=-S_{12}&=\tfrac1{12}\left(9+3\left(\frac2p\right)-\left(\frac3p\right)+\left(\frac6p\right)\right)C_1-\tfrac14C_3\\
S_2=-S_{11}&=\tfrac14\left(-1+\left(\frac2p\right)+\left(\frac3p\right)-\left(\frac6p\right)\right)C_1+\tfrac14C_3\\
S_3=-S_{10}&=\tfrac16\left(-\left(\frac3p\right)+\left(\frac6p\right)\right)C_1\\
S_4=-S_9&=\tfrac16\left(3-2\left(\frac3p\right)-\left(\frac6p\right)\right)C_1\\
S_5=-S_8&=\tfrac14\left(-1-\left(\frac2p\right)+\left(\frac3p\right)+\left(\frac6p\right)\right)C_1+\tfrac14C_3\\
S_6=-S_7&=\tfrac1{12}\left(3-3\left(\frac2p\right)+\left(\frac3p\right)-\left(\frac6p\right)\right)C_1-\tfrac14C_3
\end{aligned}
\]

If $p$ is a prime, the formulas given in the text follow from the first system; for the other cases, similar formulas are obtained from the second system.
\begin{flushright}\textsc{R. Dedekind.}\end{flushright}
\BookRule{55mm}
